\clearpage
\VerdierPhysicalPageAnchor{0120}%
\begingroup
\small
\setlength{\abovedisplayskip}{4pt plus 2pt minus 2pt}%
\setlength{\belowdisplayskip}{4pt plus 2pt minus 2pt}%

\VerdierSemanticAnchor{ch2-1-2-9-opposite-category}
$\mathrm{ii)}\Longrightarrow\mathrm{iii)}$ for the category opposite to
$\mathcal D$. The implication $\mathrm{i)}\Longrightarrow\mathrm{iii)}$
for every triangulated category is therefore equivalent to the implication
$\mathrm{ii)}\Longrightarrow\mathrm{iii)}$ for every triangulated category.
To prove the implications

\VerdierSemanticAnchor{ch2-1-2-9-reduction-implication}
\[
  \mathrm{i)}\Longleftrightarrow\mathrm{ii)}
  \Longleftrightarrow\mathrm{iii)}\Longrightarrow\mathrm{iv)}\,,
\]
it therefore suffices to prove the implication
$\mathrm{i)}\Longrightarrow\mathrm{iii)}$.

\VerdierSemanticAnchor{ch2-1-2-10}
\medskip
\noindent\textbf{1.2.10. Proof of Proposition 1.2.9 (continued).}
Let us first prove that a monomorphism is decomposable. Let
$f\colon X\longrightarrow Y$ be a monomorphism, and let
$X\xrightarrow{f}Y\xrightarrow{v}Z\xrightarrow{w}X[1]$ be a distinguished
triangle containing $f$ (TRI).

\VerdierSemanticAnchor{ch2-1-2-10-monomorphism}
Since the morphism $f$ is a monomorphism, the same is true of $f[1]$. It
then follows from TRII and (1.2.2) that $w=0$. It then follows from (1.2.6)
that the morphism $f$ is decomposable.

\VerdierSemanticAnchor{ch2-1-2-10-kernel-decomposition}
Let us now prove $\mathrm{i)}\Longrightarrow\mathrm{iii)}$ in the general
case. Suppose that the morphism $f\colon X\longrightarrow Y$ has a kernel
$\operatorname{Ker}(f)\longrightarrow X$. The morphism
$\operatorname{Ker}(f)\longrightarrow X$ is a monomorphism, hence it is
decomposable. The morphism $f$ is therefore isomorphic to a morphism of the
form

\VerdierSemanticAnchor{ch2-1-2-10-kernel-normal-form}
\[
  \operatorname{Ker}(f)\oplus X'
  \xrightarrow{(0,g)}Y\,,
\]
whose kernel is

\VerdierSemanticAnchor{ch2-1-2-10-kernel-inclusion}
\[
  \operatorname{Ker}(f)
  \xrightarrow{\left(\begin{smallmatrix}\operatorname{id}\\0\end{smallmatrix}\right)}
  \operatorname{Ker}(f)\oplus X'.
\]
It follows immediately that the morphism $g$ is a monomorphism; hence it is
decomposable. The morphism $f\colon X\longrightarrow Y$ is therefore
isomorphic to a morphism of the form

\VerdierSemanticAnchor{ch2-1-2-10-final-matrix}
\[
  \operatorname{Ker}(f)\oplus X'
  \xrightarrow{\left(\begin{smallmatrix}0&\operatorname{id}\\0&0\end{smallmatrix}\right)}
  X'\oplus Y'.
\]
It is therefore decomposable.

\VerdierSemanticAnchor{ch2-1-2-11}
\medskip
\noindent\textbf{1.2.11. End of the proof of Proposition 1.2.9.\footnote{We
use here an argument due to Freyd.}}
It remains to prove the implication
$\mathrm{iv)}\Longrightarrow\mathrm{iii)}$ under the additional
hypotheses.

\VerdierSemanticAnchor{ch2-1-2-11-duality-reduction}
First note that countable direct sums are representable in $\mathcal D$ if
and only if countable direct products are representable in
$\mathcal D^{\circ}$. Moreover, the implication
$\mathrm{iv)}\Longrightarrow\mathrm{iii)}$ for the category $\mathcal D$
is equivalent to the implication
$\mathrm{iv)}\Longrightarrow\mathrm{iii)}$ for the category opposite to
$\mathcal D$.

\VerdierSemanticAnchor{ch2-1-2-11-product-case-boundary}
It therefore suffices to prove the implication
$\mathrm{iv)}\Longrightarrow\mathrm{iii)}$ when countable products are
representable in $\mathcal D$, that is, in view of what precedes, to
\endgroup
